\subsection*{3.3. 奇相对维数}

为证明奇相对维数的 Picard--Lefschetz 公式, 我们必须归结到 XIV, 3.2.1.1 所处理的超越情形. 为此所需的特殊化论证迫使我们在不再是一维的底上工作. 我们将使用适用于所需特定情形的专门论证: 对维数大于一的底, 我并没有一套可用的理论, 所以宁愿不用泛泛之论掩饰自己的无知.

\subsubsection*{3.3.1.}
设 $S=\operatorname{Spec}(A)$ 为严格局部 Noether 概形, $f:X\to S$ 为纯相对维数为 $n$ 的有限型平坦态射, 其特殊纤维 $X_s$ 的所有非光滑点均为非退化二次点. 设 $s$ 为 $S$ 的闭点, $X_s$ 为特殊纤维, $x_0\in X_s$ 为一个奇点.

由 1.3.2(i), 在 $x_0$ 的一个邻域中, 概形 $X/S$ 局部同构于仿射空间 $\mathbb A^{n+1}_S$ 中由方程 $Q-b=0$ 定义的 $S$-概形, 其中 $b$ 属于极大理想, $Q$ 为非退化二次型. $A$ 的理想 $(b)$ 仅依赖于 $(X/S,x_0)$: 在 $x_0$ 的一个邻域中, $X$ 在 $\operatorname{Spec}(A/(b))$ 上存在一个截面
\[
s\in X(A/(b))
\]
使
\[
J^n_{X/S}\simeq
\mathcal O_{s(\operatorname{Spec}(A/(b)))}.
\]
我们将此截面记作 $s_x$, 并以 $b(x)$ 表示 $(b)$ 的任意一个生成元.

\subsubsection*{3.3.2.}
假设 $S$ 正规, 设 $b\in A$. 若 $n$ 在 $A$ 中可逆且 $b\ne0$, 则以 $\varepsilon_b^{(n)}$ 表示从 $A$ 的分式域的 Galois 群到 $\mathbb Z/(n)(1)$ 的以下同态:
\[
\varepsilon_b^{(n)}(\sigma)
= (\sigma\sqrt[n]{b})/\sqrt[n]{b}.
\]
有
\[
\bigl(\varepsilon_b^{(nm)}\bigr)^m=\varepsilon_b^{(n)}.
\]
若 $\Lambda$ 为 $\mathbb Z/n$-代数, 则仍以 $\varepsilon_b$ 表示复合
\[
\operatorname{Gal}\longrightarrow\mathbb Z/(n)(1)
\longrightarrow\Lambda(1).
\]
它与 $n$ 无关.

特征标 $\varepsilon_b$ 通过下列群分解:
\[
\pi_1\bigl(S-(\text{子概形 }b=0)\bigr).
\]

\subsubsection*{3.3.3.}
假设 $S$ 为严格 Hensel 特质, $n$ 为奇数, $n=2m+1$. 设 $x$ 为 $X_s$ 的一个奇点, 在它的一个邻域中 $X_\eta$ 光滑, 并设 $X_1/S$ 为如 1.3.2(i) 中那样由方程 $Q-b=0$ 定义的 $\mathbb A^{n+1}_S$ 的子概形. 对 $x$ 的假设意味着 $b$ 虽模极大理想为零, 本身却非零.

设 $X_{s(x)}$ 为 $X_s$ 在 $x$ 处的 Hensel 化, $\widetilde X_{s(x)}$ 为 $X_s$ 在 $x$ 处的爆破, $Y$ 为 $x$ 的逆像这一二次超曲面, 并置
\[
X^*_{s(x)}=X_{s(x)}-\{x\}=\widetilde X_{s(x)}-Y.
\]
复合态射
\begin{center}
\begin{tikzcd}[column sep=huge,row sep=huge]
H^{n-1}(Y,\Lambda(m))
  & H^{n-1}(\widetilde X_{s(x)},\Lambda(m))
      \arrow[l,"\sim"'] \arrow[r]
  & H^{n-1}(X^*_{s(x)},\Lambda(m)) \arrow[d] \\
{}
  & H^n_{\{x\}}\bigl(X_s,R\Psi_{\bar\eta}(\Lambda(m))\bigr)
  & H^n_{\{x\}}(X_s,\Lambda(m)) \arrow[l]
\end{tikzcd}
\end{center}
将后一群等同于二次超曲面 $Y$ 的上同调的本原商(XII, 3.5). 事实上, 在局部同构的情形下, 这由 2.2.7 得出. 此商的两个自然生成元($Y$ 的两族母空间的上同调类)的像为两个基向量
\[
\pm\delta\in
H^n_{\{x\}}\bigl(X_s,R\Psi_{\bar\eta}(\Lambda(m))\bigr).
\]
\footnote{定义去心 Hensel 化时, 印本在第一个 $X^*_{s(x)}$ 上加了波浪号. 它与 $X_{s(x)}-\{x\}$ 的等式以及 2.2.6 中相同的构造确定了此处使用的不带波浪号的写法.}